\documentclass[11pt,a4paper]{article}
\usepackage[utf8]{inputenc}
\usepackage[T1]{fontenc}
\usepackage[polish]{babel}
\usepackage{lmodern}
\usepackage{geometry}
\usepackage[hidelinks]{hyperref}
\geometry{margin=27mm}
\title{Zmęczenie poznawcze i sprawność uwagi u młodych dorosłych}
\author{}
\date{Prospekt badawczy, wersja 0.1.0, 22 listopada 2026}
\begin{document}
\maketitle
\begin{abstract}
Projekt MSC-FAT-P01 ma status \texttt{IDEA\_REVIEW}. Zbadać zmiany uwagi podtrzymywanej w trakcie wykonywania powtarzalnego zadania oraz ich związek z subiektywnym zmęczeniem poznawczym. Dokument zawiera planowane hipotezy i metody, a nie wyniki empiryczne.
\end{abstract}
\section{Problem i cel badania}
Zbadać zmiany uwagi podtrzymywanej w trakcie wykonywania powtarzalnego zadania oraz ich związek z subiektywnym zmęczeniem poznawczym.
\section{Hipotezy robocze}
Wraz z czasem wykonywania zadania wzrośnie liczba przeoczeń i zmienność czasu reakcji; silniejsze deklarowane zmęczenie będzie związane z większym spadkiem sprawności.
\section{Projekt badania}
Badanie eksperymentalne online dorosłych w wieku 18--30 lat; kilkublokowe zadanie uwagi podtrzymywanej typu SART albo PVT oraz krótkie oceny aktualnego zmęczenia przed zadaniem i po każdym bloku. Rejestrować typ urządzenia i porę dnia. Warunki uczestnictwa, kryteria wykluczenia, uzasadnienie liczebności próby i procedura świadomej zgody zostaną ustalone przed prerejestracją.
\section{Operacjonalizacja}
Pierwszorzędowy: zmiana liczby błędów uwagi w kolejnych blokach. Drugorzędowe: zmienność i mediana czasu reakcji, subiektywnie oceniane zmęczenie.
\section{Plan analiz}
Model mieszany dla błędów uwagi z predyktorem numeru bloku i losowym interceptem uczestnika; analiza zmęczenia jako moderatora, kontrola snu, kofeiny, pory dnia i efektów urządzeń. Przed analizą danych należy zamrozić estymand, reguły czyszczenia danych, plan braków danych i kryteria analiz wrażliwości.
\section{Etyka, bezpieczeństwo i dane}
Nie rozpoczęto rekrutacji. Konieczne jest rozstrzygnięcie wymogów opinii komisji etycznej, minimalizacji danych, retencji i informowania o dobrowolności udziału. Nie zbierać danych pozwalających na identyfikację w publicznym repozytorium.
\section{Wyniki}
Brak danych i wyników. Sekcja zostanie uzupełniona dopiero po zrealizowaniu zatwierdzonego protokołu.
\section{Ograniczenia}
Wykonanie online nie gwarantuje porównywalnych opóźnień reakcji; należy zaplanować test kalibracyjny i analizę wrażliwości. Spadek wydajności nie dowodzi sam w sobie zmęczenia.
\section{Dyskusja}
Interpretacja będzie zależała od trafności narzędzi, jakości danych i alternatywnych wyjaśnień. Projekt nie zakłada z góry potwierdzenia hipotez.
\section*{Literatura}
Weryfikacja prac źródłowych i przygotowanie bibliografii z kontrolą DOI należą do kolejnego etapu. Nie umieszczono niezweryfikowanych cytowań.
\end{document}
